\section{Introduction}
Adding vision to a language model requires choosing an encoder, an extraction layer,
and a connector. A full vision-language model (VLM) evaluation measures their combined
behavior after training. Frozen-feature probes answer a narrower question: what can a
specified readout recover at each point in the visual pipeline? This can identify
candidate representations before multimodal training, although their usefulness to a
language model must still be tested downstream.

We study DINOv2 \citep{oquab2024dinov2} and SigLIP2 \citep{tschannen2025siglip2} as
standalone references, alongside towers extracted from Kimi K2.6 \citep{moonshot2026kimik26},
Kimi K3 \citep{moonshot2026kimik3}, Muse Glimmer \citep{meta2026museglimmer}, and
Qwen3.8 \citep{qwen2026qwen38}. Recognition asks which category an image
contains. Geometry asks how far surfaces are from the camera and which direction they
face. Correspondence asks whether features match the same point across views or related
parts across objects. Occlusion tests recognition as visible image content is removed.
These capabilities need not favor the same representation.

The benchmark resolves internal layers and, for multimodal towers, the native merge and
projector stages. Towers and connectors remain frozen; recognition and geometry use
trained readouts, while correspondence uses direct feature matching. The study contributes
capability profiles for six towers, layer sweeps under two readout-capacity settings, and
paired evidence that connector effects change across evaluation datasets. The accompanying
code and result records expose the measurements and the controls behind these claims.
We report effects under this protocol rather than proposing an overall tower ranking.
